\uuid{vq1B}
\exo7id{3510}
\titre{exo7 3510}
\auteur{quercia}
\organisation{exo7}
\datecreate{2010-03-10}
\isIndication{false}
\isCorrection{true}
\chapitre{Réduction d'endomorphisme, polynôme annulateur}
\sousChapitre{Diagonalisation}
\cpge{mpsi-mp2i.matrices-applications-lineaires.a, pcsi.matrices-applications-lineaires.a, ptsi.matrices-applications-lineaires.a, mp-mpi.reduction-endomorphismes.d, pc.reduction.c, psi.reduction-endomorphismes.c, pt.reduction-endomorphismes.c, mp-mpi.reduction-endomorphismes.i, pc.reduction.d, psi.reduction-endomorphismes.d, bcpst2.valeurs-propres-vecteurs-propres.b}
\module{Algèbre}
\niveau{L2}
\difficulte{}

\contenu{
\texte{
Soit $E =  K_n[X]$ et $ u : E \to E, P \mapsto {(X^2-1)P''+(2X+1)P'.}$
}
\begin{enumerate}
    \item \question{Chercher la matrice de $u$ dans la base canonique de $ K_n[X]$.}
    \item \question{Montrer que $u$ est diagonalisable.}
\reponse{
$M = \begin{pmatrix} 0 &1 &-2 &       &(0)      \cr
                              &2 &2  &\ddots           \cr
                              &  &6  &\ddots &-n(n-1)  \cr
                              &  &   &\ddots &n        \cr
                           (0)&  &   &       &n(n+1)   \cr \end{pmatrix}$.
}
\end{enumerate}
}
